\documentclass[11pt]{article}
\usepackage[margin=1in]{geometry}
\usepackage{amsmath,amssymb,amsfonts}
\usepackage{bm}
\usepackage{booktabs}
\usepackage{enumitem}
\usepackage{xcolor}
\usepackage{hyperref}
\hypersetup{colorlinks=true,linkcolor=blue!60!black,citecolor=blue!60!black,urlcolor=blue!60!black}

\newcommand{\pderiv}[2]{\frac{\partial #1}{\partial #2}}
\newcommand{\Rav}{\widetilde{\mathrm{Ra}}}
\newcommand{\Ek}{\mathrm{Ek}}
\newcommand{\Ld}{L_d}
\newcommand{\Lc}{L_c}
\newcommand{\kh}{|\bm{k}|}
\newcommand{\Jac}[2]{J\!\left[#1,\,#2\right]}
\newcommand{\avg}[1]{\langle #1 \rangle}
\newcommand{\hhat}[1]{\widehat{#1}}
\newcommand{\Dvec}{\bm{D}}

\title{\textbf{Polar Cap Extension of the NHQGE Solver\\[4pt] via the Trap Method}\\[8pt]
\large Modifications to the Doubly-Periodic Chebyshev Solver\\for Jupiter Polar Vortex Dynamics}
\author{Formulation Notes}
\date{February 2026}

\begin{document}
\maketitle
%\tableofcontents
\vspace{1em}

%======================================================================
\section{Overview}
%======================================================================

This document describes the minimal modifications required to extend the doubly-periodic NHQGE solver (Fourier in $x,y$; Chebyshev collocation in $Z$) from a flat $\beta$-plane configuration to a polar cap geometry suitable for modeling Jupiter's polar vortex dynamics. The extension follows the ``trap'' method of Siegelman, Young \& Ingersoll (2022, PNAS), generalized from their barotropic formulation to the full three-dimensional NHQGE system.

The key idea: the polar cap geometry is encoded entirely in the background potential vorticity $\eta(x,y)$. No changes to the grid, spectral methods, time-stepping, boundary conditions, vertical discretization, or inversion are required. The solver remains doubly periodic in $(x,y)$ with Chebyshev collocation in $Z$. Dynamics are confined to a circular region within the periodic domain by a smooth PV barrier (the ``trap'').

%======================================================================
\section{From the Flat Beta-Plane to the Polar Cap}
%======================================================================

\subsection{The flat beta-plane PV equation}

In the base solver, the PV equation reads:
\begin{equation}
\pderiv{q'}{t} + \Jac{\psi}{q'} + \beta\,\pderiv{\psi}{x} - \pderiv{w}{Z} = \text{dissipation},
\label{eq:beta_plane}
\end{equation}
where $q' = \nabla_\perp^2\psi - \psi/\Ld^2$ is the perturbation PV and $\beta\psi_x$ represents advection of the background PV gradient $\beta y$ by the flow $(u,v) = (-\psi_y, \psi_x)$.

The term $\beta\psi_x$ is the special case of a \emph{linear} background PV, $\eta(x,y) = \beta y$, for which:
\begin{equation}
\Jac{\psi}{\eta} = \Jac{\psi}{\beta y} = \beta\,\psi_x.
\end{equation}

\subsection{General background PV}

For a general (nonlinear) background PV $\eta(x,y)$, the PV equation becomes:
\begin{equation}
\pderiv{q'}{t} + \Jac{\psi}{q' + \eta} - \pderiv{w}{Z} = \text{dissipation},
\label{eq:general_pv}
\end{equation}
or equivalently, since $\eta$ is time-independent:
\begin{equation}
\pderiv{q'}{t} + \Jac{\psi}{q'} + \Jac{\psi}{\eta} - \pderiv{w}{Z} = \text{dissipation}.
\label{eq:general_pv2}
\end{equation}

The only change from the flat $\beta$-plane is the replacement:
\begin{equation}
\boxed{\beta\,\psi_x \;\;\longrightarrow\;\; \Jac{\psi}{\eta} \;=\; \psi_x\,\eta_y - \psi_y\,\eta_x}
\label{eq:replacement}
\end{equation}

Everything else in the NHQGE system---the $w$ equation, the $\theta$ equation, the inversion, the vertical coupling, the boundary conditions, the dissipation---is unchanged.

\subsection{The polar cap background PV}

Near the pole of a planet with Coriolis parameter $f$ and polar radius $a_p$, the vertical component of the Coriolis parameter varies as:
\begin{equation}
f(r) = f_p\!\left(1 - \frac{r^2}{2a_p^2}\right) + O(r^4/a_p^4),
\end{equation}
where $r = \sqrt{x^2 + y^2}$ is the distance from the pole and $f_p = 2\Omega$ is the Coriolis parameter at the pole. The background (planetary) PV associated with this variation is:
\begin{equation}
\eta_{\mathrm{polar}}(r) = f(r) - f_p = -\frac{1}{2}\gamma\,r^2,
\label{eq:eta_polar}
\end{equation}
where:
\begin{equation}
\boxed{\gamma \equiv \frac{f_p}{a_p^2}}
\label{eq:gamma_def}
\end{equation}
is the single parameter characterizing the polar geometry. Note that $\gamma$ plays the role of $\beta$ but is \emph{isotropic} (radial) rather than unidirectional---it is the leading-order curvature of the Coriolis parameter near the pole.

For Jupiter: $f_p = 3.52 \times 10^{-4}\;\mathrm{s}^{-1}$, $a_p \approx 6.7\times 10^7\;\mathrm{m}$ (polar radius), giving $\gamma \approx 7.8 \times 10^{-20}\;\mathrm{m}^{-2}\,\mathrm{s}^{-1}$.

\subsection{Gradients of the polar PV}

The gradients of $\eta_{\mathrm{polar}} = -\gamma r^2/2$ are:
\begin{equation}
\eta_x = -\gamma\, x, \qquad \eta_y = -\gamma\, y.
\label{eq:eta_gradients}
\end{equation}

The Jacobian with the streamfunction is therefore:
\begin{equation}
\Jac{\psi}{\eta_{\mathrm{polar}}} = \psi_x\,\eta_y - \psi_y\,\eta_x = -\gamma\!\left(x\,\psi_y - y\,\psi_x\right) = \gamma\!\left(x\,v + y\,u\right),
\label{eq:jac_psi_eta}
\end{equation}
where $u = -\psi_y$ and $v = \psi_x$. In terms of velocity components this is $\gamma\,\bm{x}\cdot\bm{u}_\perp$, where $\bm{u}_\perp = (v, -u)$ is the velocity rotated $90^\circ$ clockwise. Physically, it represents the advection of planetary vorticity by the geostrophic flow---the polar analog of $\beta v$ on the $\beta$-plane, but with the gradient pointing radially inward toward the pole from all directions.

%======================================================================
\section{The Trap}
\label{sec:trap}
%======================================================================

\subsection{Purpose}

The doubly-periodic domain is a square $[0, L]^2$, but the polar cap dynamics should be confined to a circular region centered at the domain center. Outside this region, the dynamics should be quiescent---otherwise the periodic boundary conditions create unphysical interactions between opposite edges of the domain, which correspond to diametrically opposite points far from the pole.

Siegelman et al.\ (2022) solved this by introducing a \emph{trap}: a smooth modification of the background PV that creates a large PV step at a specified radius $r_\star$. This PV barrier acts as a dynamical wall that confines the turbulence to the interior of the trap.

\subsection{Trap construction}

Define the domain center as $(x_c, y_c) = (L/2, L/2)$ and the radial coordinate:
\begin{equation}
r = \sqrt{(x - x_c)^2 + (y - y_c)^2}.
\end{equation}

The trapped polar PV is:
\begin{equation}
\boxed{\eta(x,y) = -\frac{1}{2}\gamma\,r^2 \;\times\; \sigma_{\mathrm{trap}}(r)}
\label{eq:eta_trap}
\end{equation}
where $\sigma_{\mathrm{trap}}(r)$ is an activation function that is $\approx 1$ for $r < r_\star$ and drops smoothly to $0$ for $r > r_\star$:
\begin{equation}
\sigma_{\mathrm{trap}}(r) = \frac{1}{2}\left[1 - \tanh\!\left(A_d\,\frac{r - r_\star}{r_\star}\right)\right].
\label{eq:sigma_trap}
\end{equation}

Here:
\begin{itemize}[nosep]
\item $r_\star$ is the trap radius, typically set to $\sim 45\%$ of the domain half-width: $r_\star = 0.45 \times L/2$.
\item $A_d > 0$ controls the sharpness of the trap boundary. Siegelman et al.\ used $|A_d| = 20$ (they defined it with opposite sign convention). Larger $A_d$ gives a sharper transition.
\end{itemize}

The activation function satisfies:
\begin{equation}
\sigma_{\mathrm{trap}}(r) \approx \begin{cases}
1 & r \ll r_\star, \\
1/2 & r = r_\star, \\
0 & r \gg r_\star.
\end{cases}
\end{equation}

The resulting PV landscape has a deep ``bowl'' of negative planetary PV centered on the pole (for the North Pole, where $f_p > 0$ and cyclones are positive), with steep walls at $r = r_\star$. Cyclonic vorticity generated by convection inside the trap accumulates near the pole (due to the $\gamma$-effect driving poleward beta-drift), while the PV step at $r_\star$ prevents escape.

\subsection{Gradients of the trapped PV}

The gradients of $\eta(x,y)$ as defined in \eqref{eq:eta_trap} are:
\begin{align}
\eta_x &= -\gamma\,(x-x_c)\,\sigma_{\mathrm{trap}} - \frac{1}{2}\gamma\,r^2\,\frac{\partial\sigma_{\mathrm{trap}}}{\partial x},\label{eq:eta_x_trap}\\
\eta_y &= -\gamma\,(y-y_c)\,\sigma_{\mathrm{trap}} - \frac{1}{2}\gamma\,r^2\,\frac{\partial\sigma_{\mathrm{trap}}}{\partial y},\label{eq:eta_y_trap}
\end{align}
where:
\begin{align}
\frac{\partial\sigma_{\mathrm{trap}}}{\partial x} &= -\frac{A_d}{2r_\star}\,\mathrm{sech}^2\!\left(A_d\,\frac{r - r_\star}{r_\star}\right)\frac{x - x_c}{r},\\
\frac{\partial\sigma_{\mathrm{trap}}}{\partial y} &= -\frac{A_d}{2r_\star}\,\mathrm{sech}^2\!\left(A_d\,\frac{r - r_\star}{r_\star}\right)\frac{y - y_c}{r}.
\end{align}

\textbf{However}, in practice it is simpler and more robust to precompute $\eta(x,y)$ on the physical grid and obtain $\eta_x, \eta_y$ spectrally:
\begin{equation}
\hhat\eta_x = ik_x\,\hhat\eta, \qquad \hhat\eta_y = ik_y\,\hhat\eta.
\label{eq:eta_grad_spectral}
\end{equation}
This automatically handles the trap boundary smoothly without needing to differentiate the $\tanh$ analytically, and it applies the same spectral truncation to $\eta$ and its gradients as to the dynamical fields. The only caveat is that $\eta$ should be smooth enough (large enough $A_d^{-1}$) to be well-resolved by the horizontal grid.

\subsection{Precomputation (done once at initialization)}

The following fields are computed once on the \emph{unpadded} physical grid $(N_x \times N_x)$ and stored:
\begin{enumerate}[nosep]
\item $\eta(x_i, y_j)$ from \eqref{eq:eta_trap}.
\item $\hhat\eta(k_x, k_y) = \mathrm{rfft2}[\eta]$.
\item $\eta_x(x_i, y_j) = \mathrm{irfft2}[ik_x\,\hhat\eta]$.
\item $\eta_y(x_i, y_j) = \mathrm{irfft2}[ik_y\,\hhat\eta]$.
\end{enumerate}

Also precompute on the \emph{padded} grid $(N_x^{\mathrm{pad}} \times N_x^{\mathrm{pad}})$:
\begin{enumerate}[nosep,resume]
\item $\eta_x^{\mathrm{pad}}(x_i, y_j)$ and $\eta_y^{\mathrm{pad}}(x_i, y_j)$: the PV gradients interpolated to the dealiased physical grid. This is done by zero-padding $\hhat\eta_x$ and $\hhat\eta_y$ to $(N_x^{\mathrm{pad}}, N_x^{\mathrm{pad}}/2+1)$ and applying the padded inverse FFT, with the same normalization as the Jacobian dealiasing.
\end{enumerate}

Total storage: 2 real-valued 2D arrays on the padded grid ($\eta_x^{\mathrm{pad}}, \eta_y^{\mathrm{pad}}$). Negligible memory.

%======================================================================
\section{Modified RHS Evaluation}
\label{sec:rhs}
%======================================================================

The RHS evaluation algorithm is identical to the base solver with one modification in Step~3 (horizontal Jacobians) and the removal of the $\beta$-term from Step~4 (assembly).

\subsection{Recall: base solver RHS at each CGL point}

In the flat $\beta$-plane solver, the PV tendency at collocation point $Z_j$ in horizontal spectral space is:
\begin{equation}
R^{(q)}_j = -\hhat{F}^{(q)}_j - i\beta\,k_x\,\hhat\psi_j + (\Dvec_Z\,\bm{w})_j + \text{dissipation},
\label{eq:rhs_beta}
\end{equation}
where $\hhat{F}^{(q)}_j$ is the spectral Jacobian $\widehat{\Jac{\psi}{q'}}$ at level $j$, and the $\beta$ term is evaluated in spectral space (diagonal multiplication).

\subsection{Polar cap RHS: absorb into the Jacobian}

In the polar cap solver, replace the $\beta\psi_x$ term with $\Jac{\psi}{\eta}$. This is done by computing the Jacobian of $\psi$ with $(q' + \eta)$ instead of $q'$ alone:
\begin{equation}
R^{(q)}_j = -\widehat{\Jac{\psi_j}{q'_j + \eta}} + (\Dvec_Z\,\bm{w})_j + \text{dissipation}.
\label{eq:rhs_polar}
\end{equation}

Note that $\eta(x,y)$ is independent of $Z$---the same background PV applies at every vertical level. This is the correct QG scaling: the planetary PV gradient is barotropic.

\subsection{Implementation: two equivalent approaches}

\subsubsection{Approach A: Augment the PV field before the Jacobian}

At each CGL point $Z_j$, before computing the horizontal Jacobian, form:
\begin{equation}
\hhat{Q}_j(\bm{k}) = \hhat{q}'_j(\bm{k}) + \hhat\eta(\bm{k}).
\end{equation}
Then compute $\Jac{\psi_j}{Q_j}$ instead of $\Jac{\psi_j}{q'_j}$.

\textbf{Cost:} One complex addition per $(k_x, k_y)$ per $Z$-level. The Jacobian itself requires the same 5 FFTs as before.

\textbf{Advantage:} No additional FFTs. No additional physical-space arrays. The $\eta$ contribution to the Jacobian is automatically dealiased by the 3/2-rule since $Q$ is formed before padding.

\textbf{Pseudocode:}
\begin{verbatim}
  # At each Z_j:
  Q_hat_j = q_hat_j + eta_hat   # augment PV with background
  Jac_q_j = jacobian(psi_hat_j, Q_hat_j)   # standard dealiased Jacobian
  # No separate beta term needed
\end{verbatim}

This is the cleanest approach. It adds exactly one line of code to the existing solver.

\subsubsection{Approach B: Compute the background Jacobian separately in physical space}

During the Jacobian evaluation at $Z_j$, after transforming $\psi_x$ and $\psi_y$ to the padded physical grid, compute the additional contribution:
\begin{equation}
\Jac{\psi}{\eta}\big|_{\mathrm{phys}} = (\psi_x)^{\mathrm{pad}} \cdot \eta_y^{\mathrm{pad}} - (\psi_y)^{\mathrm{pad}} \cdot \eta_x^{\mathrm{pad}},
\end{equation}
using the precomputed $\eta_x^{\mathrm{pad}}, \eta_y^{\mathrm{pad}}$ arrays.

\textbf{Cost:} Two multiplications and one subtraction on the padded physical grid per $Z$-level. No additional FFTs (the derivatives $\psi_x, \psi_y$ are already computed for the dynamical Jacobian).

\textbf{Advantage:} The background PV contribution is computed separately, making diagnostics easier (you can output $\Jac{\psi}{q'}$ and $\Jac{\psi}{\eta}$ independently). Also avoids forming $Q = q' + \eta$ in spectral space, which at very early times when $q'$ is small could in principle introduce rounding issues (though this is not a practical concern).

\textbf{Either approach is correct.} Approach~A is recommended for simplicity.

\subsection{Summary of modifications to the RHS}

\begin{center}
\begin{tabular}{lll}
\toprule
Component & Flat $\beta$-plane & Polar cap \\
\midrule
PV Jacobian argument & $\Jac{\psi}{q'}$ & $\Jac{\psi}{q' + \eta}$ \\
Linear PV advection & $+i\beta k_x \hhat\psi_j$ (spectral) & absorbed into Jacobian \\
$w$ equation & unchanged & unchanged \\
$\theta$ equation & unchanged & unchanged \\
Inversion & unchanged & unchanged \\
Vertical derivatives & unchanged & unchanged \\
BC enforcement & unchanged & unchanged \\
Dissipation & unchanged & unchanged \\
\bottomrule
\end{tabular}
\end{center}

%======================================================================
\section{Sponge Layer (Optional Damping Outside the Trap)}
\label{sec:sponge}
%======================================================================

The PV barrier from the trap activation function provides a \emph{dynamical} confinement---vortices cannot easily cross the large PV gradient at $r = r_\star$. However, weak wave activity (Rossby waves) can still propagate outward and reflect off the periodic boundaries. To suppress this, an optional sponge layer provides explicit damping outside the trap.

\subsection{Spatially varying drag}

Add a damping term to \emph{all three} prognostic equations:
\begin{equation}
\pderiv{f}{t} = \text{RHS} - \lambda(r)\,f,
\label{eq:sponge}
\end{equation}
where $f$ is any of $q', w, \theta$, and $\lambda(r)$ is a spatially varying damping rate:
\begin{equation}
\lambda(r) = \lambda_{\max} \times \frac{1}{2}\left[1 + \tanh\!\left(A_s\,\frac{r - r_s}{r_s}\right)\right],
\label{eq:lambda_sponge}
\end{equation}
with:
\begin{itemize}[nosep]
\item $r_s > r_\star$: the sponge activation radius (just outside the trap, e.g.\ $r_s = 1.05\,r_\star$).
\item $A_s$: sharpness of sponge onset (e.g.\ $A_s = 10$).
\item $\lambda_{\max}$: maximum damping rate (e.g.\ $\lambda_{\max} \sim 10/\Delta t$ for very strong damping).
\end{itemize}

\subsection{Implementation}

The sponge is multiplicative in physical space. Precompute $\lambda(x_i, y_j)$ on the unpadded grid at initialization. At each RK4 sub-step, after computing the full RHS in spectral space, transform to physical space, subtract $\lambda \cdot f$, and transform back. Alternatively, since $\lambda(r)$ is smooth and broad-spectrum, treat the sponge implicitly:
\begin{equation}
\hhat{f}^{n+1} = \hhat{f}^* \cdot \exp\!\left(-\bar\lambda\,\Delta t\right),
\end{equation}
where $\bar\lambda$ is an effective damping rate. However, because $\lambda$ varies spatially, the implicit treatment is only approximate. The simpler approach is explicit damping within the RK4 stages.

\textbf{In practice}, the trap's PV barrier is usually sufficient and the sponge is a safety measure. Siegelman et al.\ (2022) did not report needing a sponge layer---the trap alone confined the dynamics effectively. Include the sponge as an option that can be toggled off.


%======================================================================
\section{Convective Forcing Inside the Trap}
\label{sec:forcing}
%======================================================================

The NHQGE system generates its own convective forcing through the buoyancy coupling between $w$ and $\theta$---no external stochastic forcing is needed, unlike the barotropic model of Siegelman et al.\ (2022). This is a major advantage of the 3D NHQGE over the 2D barotropic model: the convective energy injection is self-consistent and emerges from the thermal boundary conditions and the reduced Rayleigh number $\Rav$.

\subsection{Energy injection pathway}

The physical sequence is:
\begin{enumerate}[nosep]
\item The unstable temperature gradient ($\partial_Z\bar\Theta < 0$) drives convective instability at horizontal scales $\sim \Lc$ (set by $\Rav$).
\item Convection generates baroclinic kinetic energy (vertical velocity $w$, baroclinic streamfunction modes).
\item Nonlinear interactions transfer energy upscale (inverse cascade) and toward the barotropic mode.
\item The $\gamma$-effect (polar PV gradient) breaks the isotropy of the inverse cascade, driving cyclonic vorticity poleward.
\item The finite deformation radius $\Ld$ (if included) arrests the cascade at $\sim\Ld$.
\item The interplay of $\gamma$, $\Ld$, and the convective injection rate determines the emergent structure: lone polar cyclone, vortex crystal, or jets.
\end{enumerate}

\subsection{Spatial uniformity of forcing}

In the NHQGE, convection is forced by the temperature gradient $\partial_Z\bar\Theta$ which is horizontally uniform (the conduction profile $\bar\Theta = 1 - Z$). The convective forcing therefore operates uniformly across the entire periodic domain, including outside the trap. This is physically appropriate: convective energy injection occurs everywhere in Jupiter's polar weather layer, not just within the vortex crystal region.

The trap confines the \emph{barotropic} inverse cascade (which is driven by the $\gamma$-effect) without needing to confine the baroclinic convective forcing. Convection outside the trap simply generates small-scale turbulence that is damped by the sponge (if active) or remains at small scales without participating in the organized polar dynamics.

If, for physical or numerical reasons, one wishes to confine the convective forcing to the trap interior, this can be done by multiplying $\Rav$ by $\sigma_{\mathrm{trap}}(r)$---effectively turning off thermal forcing outside the trap. This is straightforward but usually unnecessary.


%======================================================================
\section{Nondimensional Parameters for the Polar Cap}
\label{sec:params_polar}
%======================================================================

\subsection{The polar parameter in NHQGE-natural units}

In the base solver, lengths are in units of $\Lc$ (convective scale), the vertical coordinate is in units of $h$ (layer depth), and time is in units of the convective turnover time. The nondimensional polar parameter is:
\begin{equation}
\gamma_{\mathrm{nd}} = \gamma_{\mathrm{dim}} \times \Lc^3 \times \frac{1}{\nu/\Lc} = \gamma_{\mathrm{dim}} \times \frac{\Lc^4}{\nu},
\end{equation}
where the factor of $\Lc^3$ comes from $[\gamma] = \mathrm{m}^{-2}\mathrm{s}^{-1}$ and we need $\gamma_{\mathrm{nd}} \times r_{\mathrm{nd}}^2$ to give an inverse-time in code units. More transparently:
\begin{equation}
\gamma_{\mathrm{nd}} = \gamma_{\mathrm{dim}} \times \frac{\Lc^2}{(2\Omega\epsilon^2)},
\end{equation}
which is the ratio of the polar PV curvature to the inverse convective time scale, at the convective length scale.

\subsection{Key length scales}

\begin{center}
\begin{tabular}{cll}
\toprule
Scale & Expression & Physical meaning \\
\midrule
$\Lc$ & (unit length in code) & Convective injection scale \\
$\Ld$ & (input parameter) & Deformation radius, cascade arrest \\
$L_\gamma = (U/\gamma)^{1/3}$ & (emergent) & Crystal radius (Siegelman et al.) \\
$r_\star$ & (input) & Trap radius \\
$L$ & (input) & Domain size \\
\bottomrule
\end{tabular}
\end{center}

The fundamental prediction of Siegelman et al.\ (2022) is that $L_\gamma$ sets the crystal radius. In the NHQGE context, $U$ (the rms velocity) is not prescribed but emerges from the convective dynamics controlled by $\Rav$ and $\sigma$. Thus the crystal radius depends on the full parameter set $(\Rav, \sigma, \gamma, \Ld)$.

\subsection{Parameter regime for Jupiter}

Using the effective (turbulent) convective scale $\Lc^{\mathrm{eff}} \sim 100\;\mathrm{km}$:
\begin{center}
\begin{tabular}{ll}
\toprule
Parameter & Value \\
\midrule
$\Rav$ & $\sim 100$ (adjustable) \\
$\sigma$ & $\sim 1$ \\
$\gamma_{\mathrm{nd}}$ & set to produce $L_\gamma \sim 50\,\Lc$ \\
$\Ld^{\mathrm{nd}}$ & $\sim 2$--$4$ (in units of $\Lc^{\mathrm{eff}}$)  \\
$r_\star$ & $\sim 0.45\,L/2$ \\
$L$ & $\gtrsim 100\,\Lc$ (to accommodate $L_\gamma$ within trap) \\
\bottomrule
\end{tabular}
\end{center}

The domain must satisfy $L \gg L_\gamma$ so that the crystal fits comfortably within the trap. A practical rule: $L \ge 4\,L_\gamma$, giving $r_\star \approx 0.45 \times L/2 \approx 0.9\,L_\gamma$.


%======================================================================
\section{Diagnostics Specific to the Polar Cap}
\label{sec:diagnostics}
%======================================================================

\subsection{Restriction to the trap interior}

All spatial diagnostics (spectra, mean profiles, integral quantities) should be restricted to the trap interior $r < r_\star$. In practice, multiply fields by $\sigma_{\mathrm{trap}}(r)$ before computing diagnostics, or use an indicator mask.

\subsection{Azimuthal decomposition}

Inside the trap, azimuthal Fourier decomposition around the pole provides physically meaningful diagnostics:
\begin{equation}
f(r, \theta) = \sum_{m=-\infty}^{\infty} \hat{f}_m(r)\,e^{im\theta},
\end{equation}
where $(r, \theta)$ are polar coordinates centered on the domain center. This is computed by interpolating the Cartesian-grid fields to a polar grid and applying an FFT in $\theta$ at each $r$.

Key diagnostics:
\begin{itemize}[nosep]
\item $m = 0$: azimuthally symmetric component (polar cyclone, mean PV profile).
\item $m = N_v$: circumpolar vortex pattern with $N_v$ vortices. Jupiter's north pole has $m = 8$, south pole $m = 5$.
\item Azimuthal energy spectrum $E(m)$: identifies the dominant crystal mode.
\item Time evolution of the $m = N_v$ amplitude: measures crystal emergence and stability.
\end{itemize}

\subsection{Vortex tracking}

Identify individual cyclones as local maxima of the barotropic vorticity field $\zeta_0 = \nabla_\perp^2\psi_0$ (depth-averaged relative vorticity). Track their positions $(x_i(t), y_i(t))$ over time to measure:
\begin{itemize}[nosep]
\item Crystal rotation rate (solid-body rotation of the vortex pattern).
\item Oscillation amplitudes (cf.\ Gavriel \& Kaspi 2022).
\item Beta-drift velocities.
\item Merger events.
\end{itemize}

\subsection{Energy spectra on the polar cap}

The isotropic energy spectrum $E(k)$ is computed as before (shell-averaging of $\kh^2|\hhat\psi|^2$), but restricted to the trap interior. An alternative is the radial energy spectrum $E(r)$ obtained by integrating $|\bm{u}|^2/2$ over annuli---this directly shows the radial structure of the crystal.

\subsection{Comparison with Siegelman et al.\ (2022) diagnostics}

For validation against the barotropic results:
\begin{itemize}[nosep]
\item Crystal radius vs.\ $L_\gamma = (U/\gamma)^{1/3}$: does the NHQGE produce the same scaling?
\item Number of vortices in the crystal: sensitive to $L_\gamma / \Ld$ ratio?
\item PV homogenization inside the crystal: does the barotropic PV tend toward uniform values between cyclones?
\end{itemize}


%======================================================================
\section{Validation Cases}
\label{sec:validation}
%======================================================================

\subsection{Case P0: Barotropic limit recovery}

Set $N_Z = 1$ (effectively barotropic), $\Ld = \infty$, prescribe an initial small-scale random vorticity field matching Siegelman et al.\ (2022) parameters. Verify that the solver reproduces their vortex crystal phenomenology: crystal radius $\sim L_\gamma$, long-lived quasi-regular array, PV homogenization.

This tests the trap implementation and the $\gamma$-effect without the complication of baroclinic dynamics.

\subsection{Case P1: Convectively forced, no polar effect}

$\gamma = 0$, $\Ld = \infty$, trap active (but irrelevant since there is no $\gamma$-driven poleward drift). This should reproduce the Rubio et al.\ (2014) isotropic condensation---a domain-filling dipolar vortex---confirming that the trap does not interfere with the inverse cascade when $\gamma = 0$.

\subsection{Case P2: Convectively forced with polar PV}

$\gamma > 0$, $\Ld = \infty$. The inverse cascade is now anisotropic. Cyclonic vorticity should accumulate near the pole (center of trap). Depending on the ratio of the convective injection scale to $L_\gamma$, expect either a lone polar cyclone (large $\gamma$) or a vortex crystal (small to moderate $\gamma$).

This is the first genuinely new result: a convectively forced polar vortex crystal, where the convective injection is self-consistent rather than prescribed.

\subsection{Case P3: Convectively forced with polar PV and finite deformation radius}

$\gamma > 0$, $\Ld$ finite. The deformation radius arrests the cascade at a scale intermediate between the convective injection and the crystal radius. This is the most Jovian-relevant configuration.

Key question: does $\Ld$ control the size of individual circumpolar cyclones (as suggested by observations showing cyclone diameters $\sim \Ld$)?

\subsection{Case P4: Parameter sweep}

Batch over $(\gamma, \Ld)$ on the H100/H200. Map the phase diagram:
\begin{itemize}[nosep]
\item Small $\gamma$, large $\Ld$: isotropic condensation (no crystal).
\item Large $\gamma$, large $\Ld$: lone polar cyclone.
\item Moderate $\gamma$, moderate $\Ld$: vortex crystal (the Jupiter regime).
\item Moderate $\gamma$, small $\Ld$: many small vortices, possibly disorganized.
\end{itemize}


%======================================================================
\section{Implementation Checklist}
\label{sec:checklist}
%======================================================================

Modifications to the existing doubly-periodic Chebyshev solver:

\begin{enumerate}[nosep]
\item \textbf{Initialization:} Compute $\eta(x,y)$ on the physical grid from \eqref{eq:eta_trap} using input parameters $\gamma$, $r_\star$, $A_d$. Compute $\hhat\eta$ via \texttt{rfft2}.

\item \textbf{RHS evaluation:} At each CGL point $Z_j$, replace
\begin{center}
\texttt{Jac\_q = jacobian(psi\_hat\_j, q\_hat\_j)} \\
$\downarrow$ \\
\texttt{Jac\_q = jacobian(psi\_hat\_j, q\_hat\_j + eta\_hat)}
\end{center}
Remove the $\beta\psi_x$ spectral term.

\item \textbf{(Optional) Sponge:} Precompute $\lambda(x,y)$ on the physical grid. Add $-\lambda\cdot f$ to the RHS of all three equations.

\item \textbf{Parameters:} Replace $\beta$ (scalar) with $\gamma$ (scalar), $r_\star$ (scalar), $A_d$ (scalar) in the configuration.

\item \textbf{Diagnostics:} Add azimuthal decomposition, vortex tracking, trap-restricted spectra.
\end{enumerate}

That's it. Five items, of which only item~2 touches the core solver, and it is a one-line change.


%======================================================================
\section{Summary}
\label{sec:summary}
%======================================================================

\noindent\fbox{\parbox{0.95\textwidth}{
\textbf{Polar cap NHQGE = base solver + one modification:}

\vspace{0.5em}
Replace the flat $\beta$-plane PV advection
\[
\beta\,\psi_x
\]
with the Jacobian against the trapped polar background PV
\[
\Jac{\psi}{\eta}, \qquad \eta(x,y) = -\tfrac{1}{2}\gamma\,r^2\,\sigma_{\mathrm{trap}}(r),
\]
implemented by augmenting $q' \to q' + \eta$ before the standard dealiased Jacobian computation.

\vspace{0.5em}
\textbf{New parameters:} $\gamma = f_p/a_p^2$ (polar PV curvature), $r_\star$ (trap radius), $A_d$ (trap sharpness).

\vspace{0.5em}
\textbf{Everything else is unchanged:} grid, FFTs, Chebyshev vertical, inversion, BCs, dissipation, time-stepping, GPU layout, batch strategy.
}}

\end{document}
